% !TeX root = ../MAgent_Theory_2.0.tex
\chapter{Communication Codes and Learned Alignment}
\label{ch:communication-links}

A communication channel can preserve a transmitted marginal while leaving the receiver to supply its conditional interpretation. An invertible dictionary instead aligns complete content spaces. Keeping these objects separate permits a precise account of lost information, learned alignment, and the distinctions that a loop of communicating agents can resolve.

The constructions take place at one context of the principal bundle $P$ of \cref{ch:foundations}. One group $G$ acts on belief and model contents through $\varrho_q$ and $\varrho_s$. A frame change is one group element acting in both sectors. General-law reconstruction uses standard Borel spaces and normalized kernels. For the matrix calculations, write the belief representation as $V=\R^K$ with $G\subseteq\GL(V)$ and suppress the sector superscript. The model-sector calculation uses its own representation and codes; using one dictionary for both sectors requires the tied-link hypothesis. Opposite attending dictionaries remain independent.

For sender $j$ and receiver $i$, selective transmission has the form
\begin{equation}
q_j\longrightarrow(T_j)_\#q_j
\longrightarrow\text{transmission noise}
\longrightarrow\text{receiver conditional law}
\longrightarrow q_i'.
\label{it1:chain}
\end{equation}
This is a law-reconstruction protocol when the message law is supplied. Conditioning on one realized message is a different operation. The general-law information identities below apply to both finite alphabets and continuous contents. Gaussian assumptions enter only the explicit decoder and tangent posterior calculations.

\section{Codes and dictionaries}
\label{it1:ch-objects}

\definitionheading{Code of a modality}{it1:code}
For general laws, a modality is a normalized Markov kernel from a standard Borel content space to a standard Borel message space. A noiseless modality is a measurable map. The linear realization used for identifiability below is more restrictive.
A modality $c$, such as speech, drawing or gesture, has a public message space $H_c=\R^{r_c}$ and a noise covariance $R_c$, either positive definite or zero for a noiseless code. The message space carries no frame of its own: its coordinates are public conventions. Sender $j$'s code in modality $c$ is a linear map $\bar T_j^c:V\to H_c$ of rank $r_c\leq K$ on reference coordinates. Its displayed form is $T_j^c=\bar T_j^cU_j^{-1}$. A message produced from displayed content $y_j$ is
\begin{equation}
h=T_j^cy_j+\varepsilon,\qquad\varepsilon\sim\mathcal N(0,R_c).
\label{it1:message}
\end{equation}

Under a passive change $U_j\mapsto g_jU_j$ the displayed content becomes $g_jy_j$ and the displayed code becomes $T_j^cg_j^{-1}$, so the message is unchanged. Under a common reference rechoice $U_j\mapsto U_jk$ with $k\in G$, the reference code becomes $\bar T_j^c k$ and the reference dictionary becomes $k^{-1}\bar\theta_{ij}k$. Re-expressing these constitutive data is the reference-consistency requirement. Public messages therefore do not depend on the sender's frame, and the reference code $\bar T_j^c$ is unchanged by passive display changes at fixed reference coordinates. Codes are agent specific. The shared generative model of \cref{sec:generative-models} contains no code, so a shared language is an additional declared hypothesis: members of a species use codes related by one group element in reference coordinates. Whether learned codes come to satisfy this hypothesis is \cref{it1:open-codes}. The identification results below state explicitly when they assume that the receiver knows the sender's reference code.

\definitionheading{Dictionary}{it1:dictionary}
Receiver $i$'s dictionary for sender $j$ is
\begin{equation}
\theta_{ij}=U_i\bar\theta_{ij}U_j^{-1}\in G,\qquad
\bar\theta_{ij}\in G.
\label{it1:dictionary-eq}
\end{equation}
This is the group-valued link of \cref{def:int-group-links} with the reference value taken in all of $G$. The general dictionary is unrestricted in $G$. The local parameterization $\bar\theta_{ij}=\exp(\omega_{ij})$ with $\omega_{ij}\in\mathfrak g$ is used for tangent learning. It can cover less than $G$: a single real matrix exponential need not cover even $\GL^+(K,\R)$ \citep{Culver1966}. The global coset and exact-identification results below concern the stated group, and a restriction to the exponential image retains only the solutions lying in that image.

It is owned by the receiver: it is $i$'s hypothesis about how $j$'s internal reference coordinates sit relative to its own. Content with reference coordinates $\bar y$ for $i$ corresponds, according to $i$, to $\bar\theta_{ij}^{-1}\bar y$ for $j$. The effective code that $i$ attributes to $j$ is therefore
\begin{equation}
\hat A_{ij}^c=\bar T_j^c\bar\theta_{ij}^{-1}.
\label{it1:effective-code}
\end{equation}

The dictionary $\theta_{ji}$ is $j$'s hypothesis about $i$. It is an independent datum, learned by a different agent, so $\theta_{ji}$ need not equal $\theta_{ij}^{-1}$. The product $\bar\theta_{ij}\bar\theta_{ji}$ is the holonomy of the two-cycle $i\to j\to i$: the composite of what $i$ thinks $j$ means with what $j$ thinks $i$ means. It is the identity exactly when the two agents' translations of each other are mutually inverse; otherwise it measures mutual misunderstanding. Geometric reversal of one particular dictionary is its inverse; the separately learned opposite attending edge is independent.

\remarkheading{Why the dictionary is group valued}{it1:why-group}
A model that $i$ holds of $j$ need not in general be a group element, and the group-valued form is a hypothesis with a stated scope. Four considerations support it within a species. First, the two agents' internal spaces are copies of one $G$-space, so a hypothesis about how the sender reads its coordinates is naturally $G$-valued; the bundle identification of the two frames is itself the value $\bar\theta_{ij}=I$, and a dictionary with nontrivial holonomy is an additional edge field rather than a relation between frames. Partial access is carried by the code, not by the dictionary. Second, one dictionary serves every modality, which is what lets several codes pool their evidence in \cref{it1:modalities}. Third, invertibility is used: the consensus theorem \cref{it1:consensus} needs inverses to show that transported beliefs are well defined along different paths, and group-valued holonomy composes invertible links; singular path products can still have trace invariants, but do not define transport by group elements. Fourth, $i$'s knowledge of $j$ is a probability law over $G$, or over $\omega_{ij}$ in a declared exponential chart; uncertainty about the dictionary is carried by that law. For agents of different kinds, whose internal spaces differ in dimension or representation, an invertible dictionary need not exist. The natural replacement is a linear or Markov map between different spaces, with consistency measured by global sections of a sheaf rather than by holonomy \citep{HansenGhrist2021}.

\section{Reception through a code}
\label{it1:ch-reception}

Fix a receiver $i$, a sender $j$ and one modality, and suppress the indices. In $i$'s reference coordinates let $p$ be $i$'s prior law of the content and let $q=(\bar\theta)_\#\bar q_j$ be the sender's belief carried into $i$'s coordinates by the dictionary. Write $A=\hat A$ for the effective code.

A conditional law decodes a message under a receiver prior. Reconstructing an entire sender law uses that decoder for every possible message. For a measurable code $T$ and laws $q,p$ on a standard Borel space, choose a regular conditional law $p(dy\mid h)$ of $p$ given $T=h$, and define the reconstruction
\begin{equation}
\hat q(dy)=\int p(dy\mid h)(T_\#q)(dh).
\label{it1:reconstruction}
\end{equation}
It keeps exactly what the code transmits, the message law $T_\#q$, and fills every discarded direction with the receiver's own conditional belief.

\corollaryheading{Meaning lost through a noiseless code}{it1:loss}
If $\KL(q\Vert p)<\infty$, then
\begin{equation}
\KL(q\Vert\hat q)=\KL(q\Vert p)-\KL(T_\#q\Vert T_\#p).
\label{it1:loss-eq}
\end{equation}
This is the conditional-lift construction of \cref{prop:conditional-lift} after transporting the sender into the receiver's reference coordinates. The decisive proof is the chain-rule argument in \cref{cor:pythagorean}: its two nonnegative terms are finite under the hypothesis here, so subtraction gives \cref{it1:loss-eq}. The code is sufficient for the pair precisely when the discarded conditional laws agree.

In words, the meaning lost when a belief is sent through a code and rebuilt by a receiver equals the full disagreement between sender and receiver minus the disagreement that the code makes visible. It vanishes exactly when the code is sufficient for the pair, that is, when the sender's conditional belief given the message already equals the receiver's. The loss is evaluated at the true dictionary, through $q=(\bar\theta)_\#\bar q_j$. It is not constant on the coset $\bar\theta\operatorname{Stab}_G(\bar T)$ of \cref{it1:identifiability} and therefore cannot be estimated from messages, whereas $\hat q$ depends on the dictionary only through the effective code.

\corollaryheading{Noisy codes}{it1:noisy}
Let a noisy code be a Markov kernel $\mathcal K$ from the content space to a standard Borel message space, for example the Gaussian kernel $\mathcal K(\bar y,\cdot)=\mathcal N(A\bar y,R)$. Let $p^h$ and $q^h$ be regular versions of the posteriors of $p$ and $q$ given $h$, and define the reconstruction $\hat q=\int p^h(q\mathcal K)(dh)$. If $\KL(q\Vert p)<\infty$, then
\begin{equation}
\KL(q\Vert\hat q)\leq\KL(q\Vert p)-\KL(q\mathcal K\Vert p\mathcal K).
\label{it1:noisy-eq}
\end{equation}
\begin{proof}
The chain rule applied to the joint laws of $(\bar y,h)$ under $q\otimes\mathcal K$ and $p\otimes\mathcal K$, which share the kernel, gives $\KL(q\Vert p)=\KL(q\mathcal K\Vert p\mathcal K)+\int\KL(q^h\Vert p^h)(q\mathcal K)(dh)$. Since $q=\int q^h(q\mathcal K)(dh)$, joint convexity of relative entropy bounds $\KL(q\Vert\hat q)$ by the integral.
\end{proof}

The proof uses only the chain rule and joint convexity, so the bound holds for every kernel and every pair of laws with $\KL(q\Vert p)<\infty$. For $q=\mathcal N(a,\Sigma_q)$, $p=\mathcal N(m,\Sigma_p)$, and the Gaussian kernel, $\hat q$ has mean $m+OA(a-m)$ and covariance $\Sigma_p-OA\Sigma_p+O(A\Sigma_qA^T+R)O^T$ with $O$ from \cref{it1:gain}.

\propositionheading{Reception of one message}{it1:kalman}
Let $p=\mathcal N(m,\Sigma_p)$ with $\Sigma_p$ positive definite, let $A$ be an $r\times K$ matrix, and let $R$ be positive definite. Given a message $h$ from the model $h=A\bar y+\varepsilon$, the posterior of $\bar y$ is $\mathcal N(m+O(h-Am),\Sigma_p-OA\Sigma_p)$ with
\begin{equation}
O=\Sigma_pA^T(A\Sigma_pA^T+R)^{-1}.
\label{it1:gain}
\end{equation}
The posterior precision is $\Sigma_p^{-1}+A^TR^{-1}A$, a correction of rank at most $r$, and the information carried by one message is
\begin{equation}
I(\bar y;h)=\tfrac12\log\det\bigl(I_r+R^{-1}A\Sigma_pA^T\bigr).
\label{it1:mutual-information}
\end{equation}
\begin{proof}
The pair $(\bar y,h)$ is jointly Gaussian with $\operatorname{Cov}(\bar y,h)=\Sigma_pA^T$ and $\operatorname{Cov}(h)=A\Sigma_pA^T+R$, and Gaussian conditioning gives the posterior \citep[Section~2.3.1]{Bishop2006}; this is the update step of the Kalman filter \citep{Kalman1960}. Since $R$ is positive definite, $A\Sigma_pA^T+R$ is invertible without a rank condition on $A$. The precision form follows from the Woodbury identity. The mutual information is the difference of the entropies of $\bar y$ and $\bar y$ given $h$, $\frac12\log\det\Sigma_p-\frac12\log\det(\Sigma_p-OA\Sigma_p)$, and the matrix determinant lemma turns this into \cref{it1:mutual-information}.
\end{proof}

The receiving map is therefore derived. It is the gain $O$, fixed by the receiver's own uncertainty $\Sigma_p$ and by its model $A$ of the sender's code. In this Gaussian instance one message changes the receiver's precision in at most $r$ directions, however large $K$ is.

This is the Gaussian instance of conditional decoding. That the decoder is a linear map, that the posterior covariance does not depend on the message received, and that one message changes the precision in at most $r$ directions are features of the Gaussian family, not of decoding in general.

\propositionheading{Gaussian reconstruction}{it1:gaussian-reconstruction}
Let $q=\mathcal N(a,\Sigma_q)$, $p=\mathcal N(m,\Sigma_p)$ with positive-definite covariances, and let $A$ have full row rank. Put
\begin{equation}
O_0=\Sigma_pA^T(A\Sigma_pA^T)^{-1},\qquad P=O_0A.
\label{it1:projector}
\end{equation}
Then $AO_0=I_r$, $P^2=P$, $\ker P=\ker A$, the range of $P$ is the range of $\Sigma_pA^T$, and the reconstruction through the noiseless code $A$ is
\begin{equation}
\hat q=\mathcal N\bigl((I-P)m+Pa,\ (I-P)\Sigma_p(I-P)^T+P\Sigma_qP^T\bigr).
\label{it1:gaussian-reconstruction-eq}
\end{equation}
\begin{proof}
The identities for $O_0$ and $P$ are immediate. Conditioning $p$ on $A\bar y=h$ gives $\mathcal N(m+O_0(h-Am),\Sigma_p-O_0A\Sigma_p)$. Mixing over $h\sim\mathcal N(Aa,A\Sigma_qA^T)$ gives mean $m+O_0A(a-m)=(I-P)m+Pa$ and covariance $\Sigma_p-P\Sigma_p+O_0A\Sigma_qA^TO_0^T$. The last term is $P\Sigma_qP^T$. The matrix $P\Sigma_p=\Sigma_pA^T(A\Sigma_pA^T)^{-1}A\Sigma_p$ is symmetric and satisfies $P\Sigma_pP^T=P\Sigma_p$, so $\Sigma_p-P\Sigma_p=(I-P)\Sigma_p(I-P)^T$.
\end{proof}

The reconstructed belief adopts the sender's belief along the communicated directions and keeps the receiver's own belief along the discarded ones. Its update of the mean,
\begin{equation}
\mu_i'=\mu_i+O_0\hat A_{ij}\bigl(\bar\theta_{ij}\bar\mu_j-\mu_i\bigr),
\label{it1:mean-update}
\end{equation}
is a projected message of the kind computed by attention. The direction in which a communicated coordinate is lifted, the range of $\Sigma_pA^T$, is set by the receiver's own correlations: two receivers with different priors read the same message as meaning different things. The linear lift is the Gaussian instance of a general-law statement: by \cref{it1:reconstruction} the receiver keeps the transmitted marginal and supplies its own conditional law given the coded statistic, so the meaning attached to a message is fixed by the receiver's conditional structure. The loss \cref{it1:loss-eq} evaluates in closed form for \cref{it1:gaussian-reconstruction-eq}.

\remarkheading{The link with a code-shaped generator}{it1:exponential}
For matrices $O$ of size $K\times r$ and $T$ of size $r\times K$, $e^{OT}=I+O\varphi_1(TO)T$ with $\varphi_1(z)=(e^z-1)/z$ \citep[Section~1.8]{Higham2008}, and $\det e^{OT}=e^{\Tr(OT)}>0$. With the noiseless gain of \cref{it1:projector}, $AO_0=I_r$ gives $e^{P}=I+(e-1)P$: the invertible link whose generator is the reconstruction projector is a dilation along the communicated directions, not the reconstruction. Reconstruction is the affine use $(I-P)m+Pa$ of the projector; as a map of the sender mean $a$ with fixed receiver mean $m$, it is noninvertible when $r<K$. For a square full-rank code, $P=I$ and this particular loss disappears. Thus no group element factors through a code of rank $r<K$, and a code-shaped generator $OT$ exponentiates to a link that acts as the identity on $\ker T$. Codes and group elements are nevertheless related directly, through the right action in $\hat A=\bar T\bar\theta^{-1}$ and through the stabilizer $\operatorname{Stab}_G(\bar T)$ of the next section.

\section{What messages reveal about the dictionary}
\label{it1:ch-identifiability}

A dictionary can be tested only against something both agents refer to. Let a referent have law $\rho$ on $V$ in $i$'s reference coordinates, with finite positive-definite second moment $M=\E_\rho[\bar y\bar y^T]$. The sender encodes the referent in its own coordinates, and the receiver models the message as
\begin{equation}
h\mid\bar y\sim\mathcal N\bigl(\bar T\bar\theta^{-1}\bar y,R\bigr),
\label{it1:grounded-model}
\end{equation}
with $R$ positive definite.

\theoremheading{Identifiability up to the code stabilizer}{it1:identifiability}
Assume that the receiver knows the sender's reference code $\bar T$, for example through a declared shared language in reference coordinates, so that only the dictionary is unknown. Let $\operatorname{Stab}_G(\bar T)=\{g\in G:\bar Tg=\bar T\}$. For parts (i) and (iii), allow any fixed additive-noise law on $\R^r$, independent of the referent; part (ii) assumes the Gaussian noise of \cref{it1:grounded-model}.

\textup{(i)} Two dictionaries $\bar\theta',\bar\theta''$ give the same conditional message law for every referent if and only if $\bar\theta'^{-1}\bar\theta''\in\operatorname{Stab}_G(\bar T)$.

\textup{(ii)} For the perturbation $\bar\theta e^{\delta}$, $\delta\in\mathfrak g$, the Fisher information of \cref{it1:grounded-model}, averaged over $\rho$, is
\begin{equation}
F(\delta,\delta)=\Tr\bigl(R^{-1}\bar T\delta C\delta^T\bar T^T\bigr),\qquad C=\bar\theta^{-1}M\bar\theta^{-T},
\label{it1:fisher}
\end{equation}
and its kernel is $\{\delta\in\mathfrak g:\bar T\delta=0\}$, the Lie algebra of $\operatorname{Stab}_G(\bar T)$.

\textup{(iii)} For $G=\GL(K,\R)$ and $\bar T$ of rank $r$, $\operatorname{Stab}_G(\bar T)$ has dimension $K(K-r)$, and the identifiable part of the dictionary has dimension $Kr$.
\begin{proof}
(i) A probability law on $\R^r$ has no nonzero translation symmetry: its characteristic function is nonzero near the origin, where translation by $a$ would require $e^{\mathrm i\langle t,a\rangle}=1$ and hence $a=0$. Thus two translates of the same noise law agree exactly when their locations agree, and $\bar T\bar\theta'^{-1}\bar y=\bar T\bar\theta''^{-1}\bar y$ for every $\bar y$ means $\bar T\bar\theta'^{-1}=\bar T\bar\theta''^{-1}$, that is, $\bar T\bar\theta'^{-1}\bar\theta''=\bar T$. The set $\operatorname{Stab}_G(\bar T)$ is a subgroup because $\bar Tg_1g_2=\bar Tg_2=\bar T$. (ii) The mean $\bar T e^{-\delta}\bar\theta^{-1}\bar y$ has derivative $-\bar T\delta\bar\theta^{-1}\bar y$ at $\delta=0$, and the Fisher information of a Gaussian mean with fixed covariance $R$ is $\E_\rho[(\bar T\delta\bar\theta^{-1}\bar y)^TR^{-1}(\bar T\delta\bar\theta^{-1}\bar y)]$, which is \cref{it1:fisher}. Since $C$ and $R^{-1}$ are positive definite, $F(\delta,\delta)=0$ exactly when $R^{-1/2}\bar T\delta C^{1/2}=0$, that is, when $\bar T\delta=0$. For such $\delta$ every power of $\delta$ is annihilated by $\bar T$, so $\bar Te^{t\delta}=\bar T$ for all $t$. (iii) $\bar Tg=\bar T$ means that $g-I$ maps $V$ into $\ker\bar T$, which has dimension $K-r$; such $g$ form an open subset of an affine space of dimension $K(K-r)$ containing the identity.
\end{proof}

Messages never test the part of the dictionary that the code discards. Relative to the observable algebra of coded messages that part is gauge in the sense of the observable-algebra convention of \cref{ch:foundations}: no message distinguishes two dictionaries that differ by an element of the stabilizer. In interpretive terms this is an instance of the indeterminacy of translation \citep{Quine1960}; the correspondence is not a theorem about that thesis. The known-code hypothesis is essential. Messages depend on the pair $(\bar T,\bar\theta)$ only through the effective code $\hat A=\bar T\bar\theta^{-1}$, and $(\bar Tg)(\bar\theta g)^{-1}=\bar T\bar\theta^{-1}$ for every $g\in G$. If the reference code is unknown, the grounded experiment therefore identifies $\hat A$ and nothing about $\bar\theta$ separately. The stated full identification uses shared referents with a positive-definite second moment. If only one referent vector were supplied, the conditional means would determine $\bar T\bar\theta^{-1}$ only on one vector. Noise rescales the Fisher information but does not change its kernel.

\remarkheading{Compact dictionary groups}{it1:compact-stabilizer}
Part (iii) gives the largest invisible part. Let $G=O(K)$ or $SO(K)$ and let $\bar T$ have rank $r$. For orthogonal $g$, $\bar Tg=\bar T$ means $g^T$, and hence $g=g^{-T}$, fixes the row space of $\bar T$ pointwise; conversely such $g$ satisfy $\bar Tg=\bar T$. Thus $\operatorname{Stab}_G(\bar T)$ is the pointwise stabilizer of $\operatorname{row}\bar T$, isomorphic to $O(K-r)$ or $SO(K-r)$ through its action on $\ker\bar T$, of dimension $(K-r)(K-r-1)/2$, and the identifiable dimension is $r(2K-r-1)/2$. A code of rank $K-1$ identifies an $SO(K)$ dictionary completely, whereas in $\GL(K,\R)$ it leaves $K$ invisible directions. The invisible, unreachable, and detectable spaces of \cref{it1:ch-learning} shrink accordingly for compact dictionary groups; their counts are computed below only for $\mathfrak{gl}(K,\R)$.

\remarkheading{Noise laws beyond the Gaussian}{it1:location}
Gaussian noise is a computational simplification here. Let the grounded model be the location family $h=\bar T\bar\theta^{-1}\bar y+\varepsilon$ with $\varepsilon$ independent of $\bar y$ and of a fixed law $\nu$ on $\R^r$. Parts (i) and (iii) of \cref{it1:identifiability} hold for every $\nu$, because a probability law on $\R^r$ differs from each of its nonzero translates: its characteristic function $\varphi$ is nonzero near the origin, and $e^{\mathrm i\langle t,a\rangle}\varphi(t)=\varphi(t)$ there forces $a=0$. If $\nu$ has a density $f$ with $\sqrt f\in H^1(\R^r)$, equivalently finite Fisher information $J_f=4\int\nabla\sqrt f(\nabla\sqrt f)^T$ (for example a continuously differentiable density with finite $J_f$, or the Laplace density), then the location family is differentiable in quadratic mean \citep[Chapter~7]{vanderVaart1998}, and so is the referent-averaged model because $M$ is finite. The matrix $J_f$ is then positive definite, since $a^TJ_fa=0$ would make $\sqrt f$ constant along $a$ and hence not square integrable. Part (ii) holds with $R^{-1}$ replaced by $J_f$,
\begin{equation}
F(\delta,\delta)=\Tr\bigl(J_f\bar T\delta C\delta^T\bar T^T\bigr),
\label{it1:fisher-location}
\end{equation}
with the same kernel, since the score of the perturbation is $\nabla\log f(\varepsilon)^T\bar T\delta\bar\theta^{-1}\bar y$ and $\varepsilon$ is independent of $\bar y$. \Cref{it1:modalities,it1:routed} and the structural results of \cref{it1:ch-learning} therefore hold for every such noise law. When $\varepsilon$ has finite covariance $R$ and integration by parts gives $\E[\nabla\log f(\varepsilon)(\varepsilon-\E\varepsilon)^T]=-I$, positive semidefiniteness of the joint covariance of the score and $\varepsilon$ gives the location Cram\'er--Rao inequality $J_f\succeq R^{-1}$, with equality exactly when the score is affine, that is, for Gaussian noise. The Gaussian value $J_f=R^{-1}$ is therefore the least Fisher information at a given covariance. For noise laws without differentiability in quadratic mean, such as uniform noise, only parts (i) and (iii) remain. Codes whose messages are not a location family, such as codes into finite alphabets, are \cref{it1:open-general}.

\section{Several modalities}
\label{it1:ch-modalities}

\corollaryheading{Pooling modalities through one dictionary}{it1:modalities}
Let modalities $c=1,\ldots,C$ have codes $\bar T^c$ and noises $R_c$, all read through the same dictionary $\bar\theta$. Assume the messages are conditionally independent given their referents and dictionary. Then the dictionary is identified up to $\bigcap_c\operatorname{Stab}_G(\bar T^c)$, and the Fisher information is $\sum_cF_c$ with kernel $\{\delta\in\mathfrak g:\bar T^c\delta=0\text{ for all }c\}$. For $G=\GL(K,\R)$ the identifiable dimension is
\begin{equation}
K\dim\Bigl(\sum_c\operatorname{row}\bar T^c\Bigr)\leq K\min\Bigl(K,\sum_cr_c\Bigr),
\label{it1:modalities-eq}
\end{equation}
If $\sum_c r_c\leq K$, equality in the bound holds exactly when the row spaces form a direct sum. If $\sum_c r_c\geq K$, equality holds exactly when their sum is $V^*$; independence of the row spaces is then unnecessary. The dictionary is fully identified exactly when the row spaces span $V^*$.
\begin{proof}
Identity of the conditional laws of all messages requires identity for each modality, which gives the intersection by \cref{it1:identifiability}. Fisher information adds over conditionally independent messages, and a sum of positive-semidefinite forms vanishes exactly where each does. For $\GL(K,\R)$, $\bar T^c\delta=0$ for all $c$ means that the range of $\delta$ lies in $\bigcap_c\ker\bar T^c$, whose dimension is $K$ minus the dimension of the sum of the row spaces.
\end{proof}

Modalities that overlap, in the sense that their row spaces share directions, add less than their widths suggest; complementary modalities, such as speech and drawing when they convey different aspects, add fully. The corollary depends on the dictionary being shared across modalities. Without it each modality would have its own effective code to learn, and the modalities could not pool their evidence about the other agent.

\Cref{it1:modalities} assumes that receiver $i$ hears sender $j$ in every modality. A receiver with several senders instead chooses, separately in each modality, whom to listen to. Let $\beta^{(c)}_{ij}$ be the probability that modality $c$ of receiver $i$ carries a message from $j$, and let $F_c^{(j)}$ denote the form \cref{it1:fisher} for the code $\bar T_j^c$, the noise $R_c$ and the dictionary $\bar\theta_{ij}$.

\corollaryheading{Routed modalities}{it1:routed}
In each exchange let a referent $\bar y\sim\rho$ be drawn, and for each modality $c$ independently let a source label $J_c$ be drawn with $\mathbb P(J_c=k)=\beta^{(c)}_{ik}$ and a message with law $\mathcal N(\bar T_k^c\bar\theta_{ik}^{-1}\bar y,R_c)$ given $J_c=k$ and $\bar y$. Consider the dictionary $\bar\theta_{ij}e^{\delta}$ of one edge $j\to i$.

\textup{(i)} If the labels are observed, the Fisher information about $\delta$ per exchange is $\sum_c\beta^{(c)}_{ij}F_c^{(j)}$.

\textup{(ii)} If the labels are latent, the Fisher information is at most $\sum_c\beta^{(c)}_{ij}F_c^{(j)}$ in the Loewner order.

\textup{(iii)} With observed labels the information vanishes exactly on $\{\delta\in\mathfrak g:\bar T_j^c\delta=0\text{ for every }c\text{ with }\beta^{(c)}_{ij}>0\}$.
\begin{proof}
(i) Given $\bar y$ and the labels, the messages of different modalities are independent, so the complete-data log likelihood is a sum over $c$. A term with $J_c=k\neq j$ does not depend on $\bar\theta_{ij}$ and contributes nothing. A term with $J_c=j$ contributes the information of \cref{it1:identifiability}. Averaging the indicator of $J_c=j$ gives the weight $\beta^{(c)}_{ij}$. (ii) The observed-data score is the conditional expectation of the complete-data score given the messages, so by the law of total variance the observed information equals the complete-data information minus the expected conditional covariance of the complete-data score given the messages, which is positive semidefinite; this is the identity of \citet{Louis1982} for incomplete-data information. (iii) A sum of positive-semidefinite forms with positive weights vanishes exactly where every form vanishes, and \cref{it1:identifiability} identifies the kernel of each.
\end{proof}

A modality that never listens to $j$ teaches $i$ nothing about its dictionary for $j$, and the directions that the attending modalities' codes discard stay invisible. When the receiver must infer who was heard, it learns less than the routing-weighted sum.

\propositionheading{A receiver's reading of a code is a dictionary modulo the stabilizer}{it1:query-dictionary}
Let $\bar T$ and $S$ be $r\times K$ matrices of full row rank $r$. Then $\{\theta\in\GL(K,\R):\bar T\theta^{-1}=S\}$ is a single left coset $\theta_0\operatorname{Stab}_{\GL(K,\R)}(\bar T)$, and for $r<K$ it meets $\GL^+(K,\R)$.
\begin{proof}
Complete the rows of $\bar T$ to a basis of the dual space and let $\mathsf C_T$ be the invertible matrix whose first $r$ rows are those of $\bar T$; define $\mathsf C_S$ likewise from $S$. With $E=(I_r\ 0)$, $\bar T=E\mathsf C_T$ and $S=E\mathsf C_S$, so $\theta_0^{-1}=\mathsf C_T^{-1}\mathsf C_S$ satisfies $\bar T\theta_0^{-1}=S$. For $r<K$, changing the sign of a completing row of $\mathsf C_S$ changes the sign of $\det\theta_0$ without changing $E\mathsf C_S$. Finally $\bar T\theta^{-1}=\bar T\theta_0^{-1}$ exactly when $\bar T\theta_0^{-1}\theta=\bar T$, that is, when $\theta_0^{-1}\theta\in\operatorname{Stab}(\bar T)$.
\end{proof}

A receiver that reads a sender's code through an arbitrary full-rank map $S$ therefore holds exactly the identifiable part of a dictionary, a point of the $rK$-dimensional quotient of \cref{it1:identifiability}. The proposition also applies to the stacked code $(\bar T^1;\ldots;\bar T^C)$ and the stacked readings $(S_1;\ldots;S_C)$. When both stacks have full row rank $\sum_cr_c\leq K$, there is one dictionary $\bar\theta$ with $S_c=\bar T^c\bar\theta^{-1}$ for every $c$, unique modulo $\bigcap_c\operatorname{Stab}(\bar T^c)$: readings learned independently for each modality are then automatically readings of one shared dictionary. For arbitrary stacks, equal ranks alone do not establish compatibility.

\propositionheading{Compatibility of stacked readings}{it1:stack-compatibility}
Let $T_*,S_*$ be $m\times K$ matrices, with no full-rank assumption. There is a dictionary $\theta\in\GL(K,\R)$ satisfying $T_*\theta^{-1}=S_*$ exactly when
\begin{equation}
\ker(T_*^T)=\ker(S_*^T).
\label{it1:stack-relations}
\end{equation}
When this holds, the solution set is $\theta_0\operatorname{Stab}_{\GL(K,\R)}(T_*)$ for any one solution $\theta_0$. If the common rank is less than $K$, the solution set meets $\GL^+(K,\R)$.
\begin{proof}
Right multiplication by an invertible matrix preserves every linear relation among rows, proving necessity. Conversely, equality of the kernels makes the map $a^TT_*\mapsto a^TS_*$ a well-defined linear isomorphism between the row spaces. Extend bases of these row spaces to bases of $(\R^K)^*$. Extending that isomorphism to the full dual space gives right multiplication by an invertible matrix $A$, with $T_*A=S_*$. Put $\theta=A^{-1}$. The coset statement follows as in \cref{it1:query-dictionary}. When the common rank is less than $K$, the extension can reverse one unused basis vector, changing the determinant sign without changing the rows of $S_*$.
\end{proof}
For example, $T_1=T_2=(1,0)$, $S_1=(1,0)$, and $S_2=(2,0)$ give two rank-one stacks but no common dictionary: the relation $T_1-T_2=0$ fails for the readings. Sharing also constrains the readings when one dictionary must serve routing and content roles, as discussed in \cref{it1:heads}.

\remarkheading{Attention heads as modalities}{it1:heads}
The separation of a sending map, a receiver reading, a decoder, and a routing row provides a structural comparison with multi-head attention \citep{Vaswani2017}. A head reads a finite-width code and routes sources, but the usual query, key, value, and output maps are independently parameterized. Such an architecture is not automatically the Bayesian message model above. Under the full-row-rank hypothesis of \cref{it1:query-dictionary}, a reading map can be expressed as a dictionary modulo its stabilizer; several readings share a dictionary precisely under \cref{it1:stack-relations}. Requiring that same dictionary for source selection and content interpretation, and deriving the decoder from the receiver prior as in \cref{it1:gain}, imposes additional restrictions. A learned output matrix alone does not establish those restrictions.

\section{Curvature as dissonance}
\label{it1:ch-dissonance}

Let the communication graph $\Gamma$ have a directed edge $j\to i$ when $i$ receives from $j$, with dictionaries $\bar\theta_{ij}$. For a directed closed walk $\gamma=(i_0\to i_1\to\cdots\to i_k=i_0)$ define its holonomy $H_\gamma=\bar\theta_{i_ki_{k-1}}\cdots\bar\theta_{i_1i_0}$, the composite of the receivers' dictionaries along the walk. Reference belief means $(\mu_i)_i$ are in consensus when
\begin{equation}
\mu_i=\bar\theta_{ij}\mu_j\qquad\text{for every edge }j\to i,
\label{it1:consensus-eq}
\end{equation}
so that every receiver, translating each sender's belief through its dictionary, finds agreement.

\theoremheading{Consensus and holonomy invariance}{it1:consensus}
Let $\Gamma$ be strongly connected and fix $i_0$. Suppose the dictionaries act as bimeasurable bijections on a standard Borel content space. The map $(q_i)_i\mapsto q_{i_0}$ is a bijection from transported consensus laws, satisfying $q_i=(\bar\theta_{ij})_\#q_j$ on every edge, onto laws invariant under every directed closed-walk holonomy based at $i_0$. In the linear realization, the map $(\mu_i)_i\mapsto\mu_{i_0}$ is a linear bijection from the solutions of \cref{it1:consensus-eq} onto $\bigcap_\gamma \operatorname{Fix}(H_\gamma)$.
\begin{proof}
Composing the edge relations around a closed walk proves that a consensus base law is invariant under its holonomy. Conversely, let $q_0$ have that invariance. For each vertex $i$, choose a directed path $p$ from $i_0$ to $i$ and set $q_i=\Pi(p)_\#q_0$, where $\Pi(p)$ is the ordered product of the dictionaries. If $p'$ is another such path, strong connectivity supplies a directed return path $s$ from $i$ to $i_0$. Both $\Pi(s)\Pi(p)$ and $\Pi(s)\Pi(p')$ are loop holonomies and preserve $q_0$. Pushforward by the bimeasurable inverse $\Pi(s)^{-1}$ gives $\Pi(p)_\#q_0=\Pi(p')_\#q_0$. The definition is independent of the path. Appending an edge to a path gives the required edge relation, and the base law determines the entire family. For vectors, replace pushforward by matrix multiplication in the same argument. This gives path-independent values $\mu_i=\Pi(p)v$ for every common fixed vector $v$, and both directions are linear.
\end{proof}

The theorem uses the invertibility of the dictionaries and the directed return paths. A two-cycle is included among its loops and constrains consensus alongside longer loops. If the matrix holonomies have no common nonzero fixed vector, the only consensus mean is zero. This does not exclude a nontrivial invariant law: a centered isotropic Gaussian, for example, is invariant under rotation holonomies. For Gaussian laws, consensus constrains both mean and covariance; for general laws, it constrains every feature moved by the holonomies. In the language of group synchronization, consensus means are the global sections of the flat bundle that the dictionaries define on $\Gamma$, and holonomy invariance is the cycle-consistency condition of that problem \citep{Singer2011,GaoBrodzkiMukherjee2021}; the directed version here keeps independent opposite dictionaries.

\corollaryheading{Frustration}{it1:frustration}
For positive weights $w_{ij}$, the consensus energy $E(\mu)=\sum_{j\to i}w_{ij}\Vert\mu_i-\bar\theta_{ij}\mu_j\Vert^2$ is a positive-semidefinite quadratic form whose kernel is the solution set of \cref{it1:consensus-eq}. If $\Gamma$ is strongly connected, its least eigenvalue on configurations of unit norm is positive exactly when the directed closed-walk holonomies based at one vertex $i_0$ have no common nonzero fixed vector.
\begin{proof}
$E(\mu)=\Vert B_\Gamma\mu\Vert^2$ for the weighted incidence operator $B_\Gamma$ with blocks $\sqrt{w_{ij}}(I,-\bar\theta_{ij})$, so $E=\mu^TB_\Gamma^TB_\Gamma\mu$ and its kernel is $\ker B_\Gamma$. The least eigenvalue of $B_\Gamma^TB_\Gamma$ is positive exactly when $\ker B_\Gamma=0$, and under strong connectivity \cref{it1:consensus} identifies $\ker B_\Gamma$ with the common fixed space of the holonomies based at $i_0$.
\end{proof}

The kernel statement holds on every graph; the holonomy criterion needs the return paths. On the acyclic graph $1\to2$, $2\to3$, $1\to3$ there is no closed walk, so the empty family of holonomies fixes every vector, while generic invertible dictionaries give $\ker B_\Gamma=0$. Under strong connectivity the holonomy groups at different base points are conjugate, so the choice of $i_0$ does not matter. For orthogonal links on undirected graphs with the inverse convention this operator is the connection Laplacian of \citet{SingerWu2012}, whose least eigenvalue measures frustration \citep{BandeiraSingerSpielman2013}. The directed version above admits independent opposite dictionaries.

\definitionheading{Dissonance of a loop}{it1:dissonance}
The dissonance of a directed closed walk $\gamma$ is the conjugacy class of $H_\gamma$. A declared elementary cycle, or plaquette in the lattice terminology of \citet{Wilson1974}, is consonant when its holonomy is the identity and dissonant otherwise. For compact $G$, a continuous conjugation-invariant defect such as $\Vert H_\gamma-I\Vert_F$ is a quantitative label that vanishes only at consonance. For $G=\GL(K,\R)$ no continuous conjugation-invariant function of $H_\gamma$ alone vanishes only at $I$: conjugating $I+E_{12}$ by $\operatorname{diag}(t,1,\ldots,1)$ and letting $t\to0$ gives $I$. The quantitative label is then the covariance-dressed defect $Q(H_\gamma;\Sigma_{i_0})$ below, with the covariance of the base agent, which is invariant under joint re-expression of the loop and the base covariance and is not a function of the conjugacy class alone. It is the loop defect of \cref{prop:int-loop-defect}, stated there for two-cycles, applied to a walk based at $i_0$; the proof is unchanged.

\remarkheading{A law comparison and a dressed link comparison}{it1:loop-comparisons}
The covariance-dressed defect
\begin{equation}
Q(H;\Sigma)=\tfrac12\Vert\Sigma^{-1/2}(I-H)\Sigma^{1/2}\Vert_F^2,\qquad \Sigma>0,
\label{it1:dressed-defect}
\end{equation}
vanishes only at $H=I$ and has a positive-definite tangent form on every matrix Lie algebra. Its nonlinear use does not introduce law-stabilizer directions. In contrast, the law-comparison penalty $\KL(q\Vert H_\#q)$ vanishes whenever $H_\#q=q$ and can have a larger null space. Its local information geometry requires a declared regular family of laws. For a quarter-turn $H$ and $q=\mathcal N(0,I_2)$, the law-comparison penalty is zero while $Q(H;I_2)=2$.

Observability depends on the record experiment. Comparing coded means tests the weaker condition $\bar T_j^c(\mu_j-\bar\theta_{ij}^{-1}\mu_i)=0$ on each edge and modality. Grounded messages with a positive-definite referent second moment instead identify the whole effective code, and hence the dictionary modulo its code stabilizer, as in \cref{it1:identifiability}. The tangent statement \cref{it1:complementary} concerns this latter experiment. A perturbation can fix one operating content while changing the code on other referents, so invisibility at one content does not imply invisibility to grounded messages. Separately, the law-comparison cost $\KL(q\Vert H_\#q)$ is blind to any holonomy stabilizing $q$. The covariance-dressed link defect \cref{it1:dressed-defect} is an observable of the link and covariance together and vanishes only at identity. These are different comparisons.

\section{Bandwidth and source rows}
\label{it1:ch-bandwidth}

\propositionheading{Coded messages give Gibbs rows at temperature one over the message count}{it1:bandwidth}
Let receiver $i$ consider finitely many candidate sources $j$ with prior $\pi_j>0$, each modeled as emitting independent messages with density $\nu_j$ on the message space. Let $h_1,h_2,\ldots$ be independent with density $\nu$, and suppose $\E_\nu|\log\nu_j|<\infty$ for every $j$. Then, almost surely,
\begin{equation}
\frac1n\log\frac{\mathbb P(J=j\mid h_{1:n})}{\mathbb P(J=k\mid h_{1:n})}\longrightarrow\E_\nu[\log\nu_j-\log\nu_k].
\label{it1:bandwidth-eq}
\end{equation}
If both $\KL(\nu\Vert\nu_j)$ and $\KL(\nu\Vert\nu_k)$ are finite, this limit is $\KL(\nu\Vert\nu_k)-\KL(\nu\Vert\nu_j)$. Integrability of $\log\nu$ under $\nu$, in addition to the candidate-log assumptions, is a sufficient condition for that representation for every pair.
\begin{proof}
By Bayes' rule the left side is $\frac1n\log(\pi_j/\pi_k)+\frac1n\sum_{t\leq n}[\log\nu_j(h_t)-\log\nu_k(h_t)]$. The strong law of large numbers sends the average to the integrable expected log-likelihood ratio in \cref{it1:bandwidth-eq}. If both relative entropies are finite, subtract their defining logarithmic integrals; the common $\log\nu$ term cancels. Without that finiteness, subtracting two infinite relative entropies is undefined.
\end{proof}

The finite-sample posterior is exactly the Gibbs row $\pi_j\exp(-nC_{j,n})/Z_n$, with empirical cost $C_{j,n}=-n^{-1}\sum_{t=1}^n\log\nu_j(h_t)$. When all candidate relative entropies are finite, its limiting relative exponential rates are those of $\pi_je^{-n\KL(\nu\Vert\nu_j)}/Z$, a Gibbs row at temperature $1/n$, the posterior concentration of a possibly misspecified model \citep{Berk1966}. The costs are divergences between message laws. When the true source and candidate $j$ emit through one common kernel $\mathcal K$, data processing gives $\KL(q\mathcal K\Vert q_j\mathcal K)\leq\KL(q\Vert q_j)$, with equality exactly when $\mathcal K$ is sufficient for the pair; when a common noiseless injective code is used, the two divergences are equal by invariance under a measurable injection. For different codes no such comparison holds: with $q=q_j=\mathcal N(0,1)$, true code $1$ and candidate code $2$, the message-law divergence is positive while the content laws agree. This identifies a cost, not a population dynamics: the normalized global model, the per-agent model, and the reciprocal potential remain distinct objectives. Finite message count and selective coding control different aspects of the source posterior.

\section{Learning the dictionaries}
\label{it1:ch-learning}

Dictionary learning requires a likelihood and a prior in addition to the communication maps. \Cref{it1:identifiability} says what messages can teach: a coset of the stabilizer of the code. This section asks what a prior that rewards agreement around loops adds to that, and answers the question exactly in the tangent model, the linearization at the identity dictionaries. Centering the prior at identity is a declared alignment hypothesis: before observing messages, the receiver expects the sender's dictionary to take its bundle-identification value $\bar\theta=I$. In interpretive terms the anchor formalizes a principle of charity in radical interpretation \citep{Davidson1973}; the correspondence is interpretive.

The tangent model is a computational simplification in two respects: it linearizes at the identity dictionaries, and it takes Gaussian noise and a Gaussian anchor. Its conclusions split accordingly. The invisible directions, the unreachable directions and their sheaf count, the split of the loop space into filled and detectable directions as a dimension count, and the exact identification results \cref{it1:exact-pair,it1:exact-flat} depend only on the identifiability kernel and on the loops; by \cref{it1:location} they hold for every additive noise law and every prior. So does the disintegration statement of \cref{it1:learning}(i): at every sample size the recognition law keeps the prior conditional law of the invisible coordinates given the identified ones. The limiting law as a conditioned Gaussian, the shrinkage and bias formulas, and the variances of \cref{it1:owned} are features of the Gaussian family.

\subsection{The tangent model}

\definitionheading{Tangent learning model}{it1:tangent-model}
Let $\Gamma$ be a finite communication graph with vertex set $\mathcal V$ and edge set $\mathcal E$, and let the edge $e:j\to i$ carry receiver $i$'s generator $\omega_e=\omega_{ij}\in\mathfrak g\subseteq\mathfrak{gl}(K,\R)$, so that $\omega=(\omega_e)_e\in X=\mathfrak g^{\mathcal E}$. Write $\bar T_e$ for the stacked code of the modalities with $\beta^{(c)}_{ij}>0$ in \cref{it1:routed}. Let $\mathcal L$ be a finite set of directed closed walks of $\Gamma$ that contains every two-cycle, each walk listed once up to the cyclic choice of its starting vertex, so that $i\to j\to i$ and $j\to i\to j$ are one element of $\mathcal L$, with one form $W_\gamma$ and, when that form is dressed by a covariance, one declared base agent. Define the linearized holonomy $B:X\to\mathfrak g^{\mathcal L}$ by $(B\omega)_\gamma=\sum_{e\in\gamma}\omega_e$, an edge being counted as often as $\gamma$ traverses it. The tangent learning model consists of the prior law with density proportional to
\begin{equation}
\exp\Bigl(-\tfrac12\omega^T\Pi_0\omega-\tfrac\lambda2\sum_{\gamma\in\mathcal L}\Vert(B\omega)_\gamma\Vert_{W_\gamma}^2\Bigr),
\label{it1:loop-prior}
\end{equation}
where $\Pi_0$ is positive definite on $X$, $\lambda\geq0$, and each $W_\gamma$ is a positive-definite quadratic form on $\mathfrak g$, with $W=\bigoplus_\gamma W_\gamma$; and of the messages that $i$ attributes to $j$, with observed labels as in \cref{it1:routed}(i), each produced from a referent $\bar y\sim\rho$ with finite second moment $M$ positive definite as
\begin{equation}
h=\bar T_e^c(I-\omega_e)\bar y+\varepsilon,\qquad\varepsilon\sim\mathcal N(0,R_c).
\label{it1:tangent-message}
\end{equation}

The prior is Gaussian with precision $\Lambda_0=\Pi_0+\lambda B^TWB$, which is positive definite, so the prior is normalized for every $\lambda\geq0$. The message model is the first-order form of \cref{it1:grounded-model} at $\bar\theta_e=e^{\omega_e}$, and $(B\omega)_\gamma$ is the first-order form of $\log H_\gamma$. By the Baker--Campbell--Hausdorff formula both errors are of second order in $\omega$. Inside the tangent model the likelihood is Gaussian and linear in $\omega$. The recognition law, here the exact posterior, is therefore Gaussian after any number of messages and is updated one message at a time by the information filter
\begin{equation}
\Lambda\leftarrow\Lambda+J_t^TR_c^{-1}J_t,\qquad
\eta\leftarrow\eta+J_t^TR_c^{-1}\bigl(h_t-\bar T_e^c\bar y_t\bigr),\qquad
J_t\omega=-\bar T_e^c\omega_e\bar y_t,
\label{it1:information-filter}
\end{equation}
started at $(\Lambda_0,0)$, with mean $\Lambda^{-1}\eta$ and covariance $\Lambda^{-1}$. This is one normalized joint over the dictionaries of all receivers, so the tangent learning model is the global picture of \cref{ch:dynamics}. The receiver-owned restriction of its recognition law is \cref{it1:owned}; per-agent models in which peers' laws are data are a different objective, treated in \cref{ch:link-dynamics}.

The loop prior is a Gaussian, noncompact lattice gauge action in the linearized holonomies \citep{Wilson1974} and a cycle-consistency penalty in the sense of group synchronization \citep{GaoBrodzkiMukherjee2021}. Its anchor removes the zero modes $\ker B$ that the loop term alone leaves free. The kernels, ranks and dimension counts below do not depend on how often a walk is listed in $\mathcal L$, whereas the shrinkage factors, the bias and the detectable dissonance do; \cref{it1:pair} lists the two-cycle once.

A change of the common reference coordinates by $k\in G$ acts by $\omega_e\mapsto k^{-1}\omega_ek$ under reference consistency, and then $(B\omega)_\gamma\mapsto k^{-1}(B\omega)_\gamma k$. The loop term is invariant when each $W_\gamma$ is. The Frobenius norm is invariant under orthogonal $k$. For $\GL(K,\R)$ the dressed norm $\Vert X\Vert_\Sigma^2=\Tr(X^T\Sigma^{-1}X\Sigma)$, with $\Sigma$ the reference covariance of the belief of the agent at which $\gamma$ is based, satisfies $\Vert k^{-1}Xk\Vert^2_{k^{-1}\Sigma k^{-T}}=\Vert X\Vert^2_\Sigma$. A law is needed here because $\mathfrak{gl}(K,\R)$ carries no conjugation-invariant positive-definite form for $K\geq2$: conjugation by $\operatorname{diag}(t,1,\ldots,1)$ multiplies $E_{12}$ by $t$. In this model $\Sigma$ is fixed data at the operating point. If it were a variable of an objective, the log normalizer $\frac12\log\det(\Pi_0+\lambda B^TW(\Sigma)B)$ of the loop prior would be a belief-dependent term that must be retained.

The anchor $\Pi_0$ and the forms $W_\gamma$ are declared in the common reference presentation of \cref{ch:foundations}; reference coordinates are common to all agents, not owned by the receiver. A fixed non-invariant anchor, such as $\sigma^{-2}I$, is therefore a choice of that presentation modulo the subgroup preserving it, which is $O(K)$ for a Frobenius form. The quantitative conclusions below, namely \cref{it1:learning}(iii)--(v), \cref{it1:complementary}(ii) and \cref{it1:pair}, hold in the presentation in which the anchor and loop forms take their displayed values and are invariant only under reference changes preserving those forms. With the covariance-dressed form used for both the anchor and the loop term they are covariant under every reference change. The kernels, ranks and zero sets are invariant in either case.

For the statements below, let $N_e=\{\delta\in\mathfrak g:\bar T_e\delta=0\}$, the Lie algebra of the stabilizer of the stacked code, and let $N=\bigoplus_eN_e\subseteq X$, the invisible directions. Let $Z$ be a matrix whose columns form an orthonormal basis of $N$ for the Frobenius inner product on $X$, and $Y$ one for $N^\perp$. Put
\begin{equation}
S=Z^T\Lambda_0Z,\qquad
\Pi_N=ZS^{-1}Z^T\Lambda_0,\qquad
L_N=Z^TB^TWBZ,
\label{it1:learning-notation}
\end{equation}
so that $\Pi_N$ is the $\Lambda_0$-orthogonal projection of $X$ onto $N$ and $L_N$ is the loop operator on the invisible directions. Let $\bar F_e=\sum_c\beta^{(c)}_{ij}F_c^{(j)}$ be the information per exchange of \cref{it1:routed}(i), evaluated at $\bar\theta_e=I$, where $C=M$.

\subsection{Concentration and selection}

\theoremheading{Learning with a loop prior}{it1:learning}
In the tangent learning model let the truth be $\omega^\star\in X$, let each edge $e$ receive $n_e=\lceil\upsilon_en\rceil$ independent exchanges with sampling rate $\upsilon_e>0$, and write $z^\star=Z^T\omega^\star$.

\textup{(i)} For every set of messages the likelihood is constant on each coset $\omega+N$, so the density of the recognition law relative to the prior depends on $\omega$ only through $Y^T\omega$. Consequently, for every prior and every additive noise law for which the messages have densities, the conditional law of $Z^T\omega$ given $Y^T\omega$ is the same under the recognition law as under the prior.

\textup{(ii)} Almost surely, as $n\to\infty$, the mean of the recognition law converges to
\begin{equation}
\omega_\infty=(I-\Pi_N)\omega^\star,
\label{it1:limit-mean}
\end{equation}
its covariance is $ZS^{-1}Z^T+O(1/n)$, and $n$ times its covariance in the coordinates $Y^T\omega$ converges to $(Y^TF_\upsilon Y)^{-1}$ with $F_\upsilon=\bigoplus_e\upsilon_e\bar F_e$. Equivalently, the recognition law converges to the prior conditioned on the coset $\omega^\star+N$. Its precision grows at rate $n$ in the identified directions, and within the coset it is the loop prior restricted to the coset.

\textup{(iii)} Let $\Pi_0=\sigma^{-2}I$. In the coordinates $Z^T\omega$ the limiting covariance is $\sigma^2(I+\lambda\sigma^2L_N)^{-1}$. It equals $\sigma^2$ on $\ker L_N=Z^T(N\cap\ker B)$, which is the prior covariance there. For $\lambda>0$ it is strictly smaller than $\sigma^2$ on the orthogonal complement, with eigenvalues $\sigma^2/(1+\lambda\sigma^2\mu)$ for the eigenvalues $\mu>0$ of $L_N$, and it tends to zero there as $\lambda\to\infty$. It never exceeds the prior marginal covariance $Z^T\Lambda_0^{-1}Z$ of the invisible coordinates and can equal it. The rank of $L_N$ is $\dim B(N)=\dim N-\dim(N\cap\ker B)$.

\textup{(iv)} If in addition the truth is flat, $B\omega^\star=0$, then
\begin{equation}
Z^T(\omega_\infty-\omega^\star)=-(I+\lambda\sigma^2L_N)^{-1}z^\star.
\label{it1:shrinkage}
\end{equation}
The loop term recovers the invisible part of the truth along an eigenvector of $L_N$ with eigenvalue $\mu>0$ up to the factor $\lambda\sigma^2\mu/(1+\lambda\sigma^2\mu)$, and as $\lambda\to\infty$ it recovers all of it except the component in $\ker L_N$, which stays at the prior mean zero.

\textup{(v)} For a general truth, as $\lambda\to\infty$,
\begin{equation}
Z^T(\omega_\infty-\omega^\star)\longrightarrow-P_{\ker}z^\star-L_N^+Z^TB^TWB\omega^\star,
\label{it1:bias}
\end{equation}
with $P_{\ker}$ the orthogonal projection onto $\ker L_N$ and $L_N^+$ the pseudoinverse.

\textup{(vi)} For every positive-definite anchor $\Pi_0$, as $\lambda\to\infty$ the coordinate $Z^T(\omega_\infty-\omega^\star)$ converges to the unique minimizer of $\Vert\omega^\star+Zz\Vert_{\Pi_0}^2$ over the minimizers of $\Vert B(\omega^\star+Zz)\Vert_W^2$. If the truth is flat and $N\cap\ker B=0$, then $\omega_\infty\to\omega^\star$.
\begin{proof}
(i) By \cref{it1:tangent-message} the law of a message on edge $e$ depends on $\omega$ only through $\bar T_e^c\omega_e$ for attending modalities $c$, which does not change when $\omega_e$ changes by an element of $N_e$. Bayes' rule gives the second statement.

(ii) After the messages the precision is $\Lambda_n=\Lambda_0+nF_n$ with $nF_n=\sum_tJ_t^TR_c^{-1}J_t$, and since $h_t-\bar T_e^c\bar y_t=J_t\omega^\star+\varepsilon_t$, the information vector is $\eta_n=nF_n\omega^\star+\xi_n$ with $\xi_n=\sum_tJ_t^TR_c^{-1}\varepsilon_t$. Every $J_t$ vanishes on $N$, so $F_n=YF_{Y,n}Y^T$ and $Z^T\xi_n=0$. By the strong law of large numbers $F_{Y,n}\to Y^TF_\upsilon Y$ almost surely, and this limit is positive definite because the kernel of $F_\upsilon$ is $N$ by \cref{it1:routed}(iii). In the basis $(Y,Z)$ the matrix $\Lambda_n$ has blocks $A_n=Y^T\Lambda_0Y+nF_{Y,n}$, $D=Y^T\Lambda_0Z$ and $S$, and block inversion gives
\begin{equation}
\Lambda_n^{-1}=ZS^{-1}Z^T+(Y-ZS^{-1}D^T)(A_n-DS^{-1}D^T)^{-1}(Y-ZS^{-1}D^T)^T,
\end{equation}
with $n(A_n-DS^{-1}D^T)^{-1}\to(Y^TF_\upsilon Y)^{-1}$. This proves the statements about the covariance, since $Y^T\Lambda_n^{-1}Y=(A_n-DS^{-1}D^T)^{-1}$. For the mean, $nF_n\omega^\star=\Lambda_n\omega^\star-\Lambda_0\omega^\star$ gives
\begin{equation}
\Lambda_n^{-1}\eta_n=\omega^\star-\Lambda_n^{-1}\Lambda_0\omega^\star+\Lambda_n^{-1}\xi_n.
\end{equation}
The second term converges to $ZS^{-1}Z^T\Lambda_0\omega^\star=\Pi_N\omega^\star$. Because $Z^T\xi_n=0$, the third term equals $(Y-ZS^{-1}D^T)(A_n-DS^{-1}D^T)^{-1}Y^T\xi_n$, which tends to zero since $\xi_n/n\to0$ almost surely by the strong law. Finally, the Gaussian law with precision $\Lambda_0$ conditioned on the affine set $\omega^\star+N$, in the coordinate $z$ of $\omega^\star+Zz$, has precision $S$ and the mean that minimizes $\frac12(\omega^\star+Zz)^T\Lambda_0(\omega^\star+Zz)$, which is $z=-S^{-1}Z^T\Lambda_0\omega^\star$. Its mean is $\omega_\infty$ and its covariance is $ZS^{-1}Z^T$.

(iii) Here $S=\sigma^{-2}I+\lambda L_N$, and the eigenvalues of $S^{-1}$ are $\sigma^2/(1+\lambda\sigma^2\mu)$ for the eigenvalues $\mu\geq0$ of $L_N$. Since $W$ is positive definite, $L_Nz=0$ exactly when $BZz=0$, so $\ker L_N=Z^T(N\cap\ker B)$ and the rank of $L_N$ is that of $BZ$.

(iv) From (ii) and $Z^T\Lambda_0\omega^\star=\sigma^{-2}z^\star+\lambda Z^TB^TWB\omega^\star$,
\begin{equation}
Z^T(\omega_\infty-\omega^\star)=-S^{-1}\bigl(\sigma^{-2}z^\star+\lambda Z^TB^TWB\omega^\star\bigr),
\label{it1:limit-error}
\end{equation}
which is \cref{it1:shrinkage} when $B\omega^\star=0$.

(v) The vector $Z^TB^TWB\omega^\star$ lies in the range of $(W^{1/2}BZ)^T$, which is the range of $L_N=(W^{1/2}BZ)^T(W^{1/2}BZ)$. As $\lambda\to\infty$, $\lambda S^{-1}\to L_N^+$ and $S^{-1}\to0$ on the range of $L_N$, while $\sigma^{-2}S^{-1}=I$ on $\ker L_N$. Applying this to \cref{it1:limit-error} gives \cref{it1:bias}.

(vi) By (ii), $z_\lambda=Z^T(\omega_\infty-\omega^\star)$ minimizes $f_\lambda=g+\lambda h$ with $g(z)=\frac12\Vert\omega^\star+Zz\Vert_{\Pi_0}^2$ and $h(z)=\frac12\Vert B(\omega^\star+Zz)\Vert_W^2$. The set $\mathcal A$ of minimizers of $h$ is a nonempty affine set, and $g$ is strictly convex and coercive, so it has a unique minimizer $z_{\mathcal A}$ on $\mathcal A$. From $f_\lambda(z_\lambda)\leq f_\lambda(z_{\mathcal A})$ one obtains $g(z_\lambda)\leq g(z_{\mathcal A})$ and $\lambda(h(z_\lambda)-\min h)\leq g(z_{\mathcal A})$. Hence $z_\lambda$ is bounded, $h(z_\lambda)\to\min h$, and every limit point lies in $\mathcal A$ with $g$ at most $g(z_{\mathcal A})$, so it equals $z_{\mathcal A}$. For a flat truth $\min h=0$ is attained at $z=0$, and $N\cap\ker B=0$ makes $BZ$ injective, so $\mathcal A=\{0\}$.
\end{proof}

Parts (i) and (ii) are the linear Gaussian instance of the limiting posterior of a partially identified model: the identified coordinates concentrate, and the nonidentified ones keep their conditional prior \citep{Poirier1998,Gustafson2010}; the underlying conditional-independence statement is that of \citet{Dawid1979}. The loop prior acts only through that conditional prior, which (iii) to (vi) compute. For a general prior and location noise the same disintegration holds at every sample size by (i); its limit requires posterior consistency of the identified coordinates and continuity of the prior conditional, which are not proved here.

With an isotropic anchor and no loop term, the recognition law on the invisible coordinates stays at its anchor marginal. For a general anchor, messages can condition these coordinates through prior correlations, while leaving their conditional law given the identified coordinates unchanged. The loop term changes the geometry in which the identified part of the truth is completed, from the anchor to the anchor plus the loop energy, and \cref{it1:limit-mean} is the resulting projection.

\subsection{Filling and detection}

\definitionheading{Detectable dissonance}{it1:detectable}
The detectable dissonance of a configuration $\omega$ is
\begin{equation}
\mathcal D(\omega)=\min_{\nu\in N}\Vert B(\omega+\nu)\Vert_W^2,
\label{it1:detectable-eq}
\end{equation}
the least linearized loop energy among the configurations that no message distinguishes from $\omega$.

\propositionheading{Filling and detection are complementary}{it1:complementary}
In the tangent learning model let $Q_{\mathcal D}$ be the $W$-orthogonal projection of $\mathfrak g^{\mathcal L}$ onto the $W$-orthogonal complement of $B(N)$.

\textup{(i)} $\mathcal D(\omega)=\Vert Q_{\mathcal D}B\omega\Vert_W^2$. In particular $\mathcal D$ is constant on cosets of $N$, so messages determine it.

\textup{(ii)} For every positive-definite anchor, as $\lambda\to\infty$ the limit mean $\omega_\infty$ satisfies $B\omega_\infty\to Q_{\mathcal D}B\omega^\star$, so its loop energy tends to $\mathcal D(\omega^\star)$.

\textup{(iii)} The loop space $B(X)$ is the $W$-orthogonal sum of $B(N)$, of dimension $\dim N-\dim(N\cap\ker B)$, and of $B(X)\cap B(N)^{\perp_W}$, of dimension $\operatorname{rank}B-\dim B(N)$. The loop energy splits as $\Vert B\omega\Vert_W^2=\Vert(I-Q_{\mathcal D})B\omega\Vert_W^2+\Vert Q_{\mathcal D}B\omega\Vert_W^2$. The first term is the only one that depends on the invisible coordinates, and a dissonance in $B(N)$ is not detectable; the second is a function of the identified coordinates alone.
\begin{proof}
(i) Since $B(\omega+Zz)=B\omega+BZz$, the minimum over $z$ leaves the $W$-orthogonal residual of $B\omega$ from the range of $BZ$, which is $B(N)$; this residual is $Q_{\mathcal D}B\omega$. For $\nu\in N$, $Q_{\mathcal D}B\nu=0$. (ii) By \cref{it1:learning}(vi), $Z^T(\omega_\infty-\omega^\star)$ converges to a minimizer $z$ of $\Vert B\omega^\star+BZz\Vert_W^2$. Every such minimizer satisfies $BZz=-(I-Q_{\mathcal D})B\omega^\star$, because $I-Q_{\mathcal D}$ is the $W$-orthogonal projection onto the range of $BZ$. Hence $B\omega_\infty\to Q_{\mathcal D}B\omega^\star$. For the isotropic anchor the same limit follows from \cref{it1:bias}. (iii) The dimension of $B(N)$ is the rank of $BZ$, computed in \cref{it1:learning}(iii). The energy splits because $Q_{\mathcal D}$ is a $W$-orthogonal projection, and $Q_{\mathcal D}B\omega$ does not depend on the component of $\omega$ in $N$ by (i).
\end{proof}

The proposition states the cost of the loop prior. Every loop direction in which agreement around loops completes a dictionary is a direction in which the community assumes an agreement it cannot check. If the truth is dissonant there, a strong loop prior can select a flat representative of the same message laws, as quantified by the bias term of \cref{it1:bias}. Dissonance is detectable only in the loop directions that the invisible directions cannot reach.

\subsection{Which directions loops cannot reach}

\propositionheading{The unreachable directions}{it1:unreachable}
Let $\Gamma$ be connected, and suppose that the chains of the walks in $\mathcal L$ span the cycle space of $\Gamma$, regarded as an undirected multigraph in which the edge $e:j\to i$ is oriented from $j$ to $i$. Then
\begin{equation}
\ker B=\bigl\{(a_i-a_j)_{e:j\to i}:a\in\mathfrak g^{\mathcal V}\bigr\},
\label{it1:coboundaries}
\end{equation}
and $N\cap\ker B$ consists of the configurations in \cref{it1:coboundaries} with $\bar T_e(a_i-a_j)=0$ on every edge $e:j\to i$. For $\mathfrak g=\mathfrak{gl}(K,\R)$,
\begin{equation}
\begin{aligned}
\dim(N\cap\ker B)&=K\bigl(\dim H^0-K\bigr),\\
H^0&=\bigl\{u\in(\R^K)^{\mathcal V}:\bar T_e(u_i-u_j)=0\text{ for every edge }e:j\to i\bigr\}.
\end{aligned}
\label{it1:sheaf-count}
\end{equation}
\begin{proof}
The value $(B\omega)_\gamma$ is the pairing of $\omega$, read as a $\mathfrak g$-valued cochain on edges, with the chain of $\gamma$, and that chain is a cycle because a closed walk leaves each vertex as often as it enters it. If these chains span the cycle space, $\omega\in\ker B$ exactly when $\omega$ vanishes on every cycle. Choose a spanning tree and a root, let $a_v$ be the sum of $\omega$ along the tree path from the root to $v$, each edge taken with the sign of its orientation, and put $a$ equal to zero at the root. Then $\omega_e-(a_i-a_j)$ vanishes on tree edges and, on any other edge, equals the value of $\omega$ on its fundamental cycle, which is zero. Conversely every configuration of the form \cref{it1:coboundaries} vanishes on cycles. The condition for membership in $N$ is the definition of $N_e$. For $\mathfrak{gl}(K,\R)$ the condition $\bar T_e(a_i-a_j)=0$ holds column by column, so $a\mapsto(a_i-a_j)_e$ maps the families whose $K$ columns lie in $H^0$, a space of dimension $K\dim H^0$, onto $N\cap\ker B$. Its kernel consists of the constant families, of dimension $K^2$, because $\Gamma$ is connected.
\end{proof}

The configurations in \cref{it1:coboundaries} are the coboundary, or vertex-wise, configurations: what the dictionaries would become if each agent's reading were re-expressed by $e^{a_v}$ with codes and contents re-read accordingly. Relative to the observable algebra of \cref{ch:foundations} they are not gauge. Passive changes leave reference dictionaries unchanged, and grounded messages distinguish a coboundary configuration from the identity unless the codes absorb it. The space $N\cap\ker B$ consists of the coboundary configurations that no attending code registers. They create no holonomy, so no loop term reaches them, and no message reveals them. \Cref{it1:coboundaries} is the coboundary part of the combinatorial Hodge decomposition of edge cochains with Lie-algebra coefficients \citep{Lim2020}. The space $H^0$ is the space of global sections of a cellular sheaf on $\Gamma$ whose stalk at an agent is $\R^K$, whose stalk at an edge is the message space, and whose two restriction maps at an edge are both the attending code \citep{HansenGhrist2019}. Such sheaves are used for discourse by \citet{HansenGhrist2021}; here the code sheaf arises from identifiability.

\corollaryheading{Complete exchange}{it1:complete-exchange}
Let every pair of agents exchange in both directions, and let $\mathcal L$ contain every two-cycle and every directed triangle. Then the hypothesis of \cref{it1:unreachable} holds and
\begin{equation}
H^0=\bigl\{u:u_i-u_j\in\ker\bar T_{j\to i}\cap\ker\bar T_{i\to j}\text{ for every pair }\{i,j\}\bigr\},
\end{equation}
where $\bar T_{j\to i}$ is the code on the edge $j\to i$. If the pairs with $\ker\bar T_{j\to i}\cap\ker\bar T_{i\to j}=0$ connect all agents, then $H^0$ consists of constant sections, $N\cap\ker B=0$, and for a flat truth in the Gaussian tangent model with any positive-definite anchor the recognition law recovers every dictionary in the limits $n\to\infty$ and then $\lambda\to\infty$.
\begin{proof}
Fix an agent $0$. The edges $0\to i$ form a spanning tree. The fundamental cycle of $i\to0$ is the two-cycle through $0$ and $i$, and that of $j\to i$ with $i,j\neq0$ is the directed triangle $0\to j\to i\to0$ minus the two-cycle through $0$ and $i$. The two edges between $i$ and $j$ give $\bar T_{j\to i}(u_i-u_j)=0$ and $\bar T_{i\to j}(u_j-u_i)=0$. If the pairs with trivial intersection connect the agents, $u$ is constant along each of them, hence constant. The last statement is \cref{it1:learning}(vi).
\end{proof}

The sufficient condition is not necessary. For three agents in complete exchange, with every two-cycle and both directed triangles in $\mathcal L$, write $I_{ij}=\ker\bar T_{j\to i}\cap\ker\bar T_{i\to j}$. A section is fixed by $u_0$ and by differences $d_{01}\in I_{01}$, $d_{12}\in I_{12}$, $d_{20}\in I_{20}$ with zero sum, so the kernel of the sum map gives
\begin{equation}
\dim(N\cap\ker B)=K\bigl[\dim I_{01}+\dim I_{12}+\dim I_{20}-\dim(I_{01}+I_{12}+I_{20})\bigr].
\label{it1:three-agent-count}
\end{equation}
In $K=3$, three linearly independent lines $I_{ij}$ give zero, although no pairwise intersection is zero: triangles remove invisible flat directions that remain for each pair separately. Three distinct coplanar lines leave $K$ such directions, and a common unread line leaves $2K$.

\propositionheading{Reciprocal pair}{it1:pair}
Let two agents exchange in both directions with $\mathcal L$ the two-cycle, $W$ the Frobenius norm and $\mathfrak g=\mathfrak{gl}(K,\R)$. Write $\bar T_1$ for the code on $1\to0$ and $\bar T_0$ for the code on $0\to1$, $\mathcal K_a=\ker\bar T_a$, $\mathcal R_a=\mathcal K_a^\perp$ for the row space of $\bar T_a$, and $P_\cap$ for the orthogonal projection onto $\mathcal R_0\cap\mathcal R_1$. Then
\begin{equation}
\begin{aligned}
\dim B(N)&=K\dim(\mathcal K_0+\mathcal K_1),\qquad
\dim(N\cap\ker B)=K\dim(\mathcal K_0\cap\mathcal K_1),\\
\operatorname{rank}B-\dim B(N)&=K\dim(\mathcal R_0\cap\mathcal R_1),
\end{aligned}
\label{it1:pair-split}
\end{equation}
and $\mathcal D(\omega)=\Vert P_\cap(\omega_{01}+\omega_{10})\Vert_F^2$.
\begin{proof}
$N$ consists of the pairs $(\nu_{01},\nu_{10})$ with the range of $\nu_{01}$ in $\mathcal K_1$ and the range of $\nu_{10}$ in $\mathcal K_0$, and $B\nu=\nu_{01}+\nu_{10}$. Hence $B(N)$ is the set of matrices with range in $\mathcal K_0+\mathcal K_1$, of dimension $K\dim(\mathcal K_0+\mathcal K_1)$. Its Frobenius-orthogonal complement in $B(X)=\mathfrak{gl}(K,\R)$ is the set of matrices with range in $(\mathcal K_0+\mathcal K_1)^\perp=\mathcal R_0\cap\mathcal R_1$, and the orthogonal projection onto it is $X\mapsto P_\cap X$. Thus $Q_{\mathcal D}=P_\cap$ and \cref{it1:complementary}(i) gives $\mathcal D$. The space $N\cap\ker B$ consists of the pairs with $\nu_{01}=-\nu_{10}$ and range in $\mathcal K_0\cap\mathcal K_1$.
\end{proof}

A pair detects its mutual misunderstanding only in $\mathcal R_0\cap\mathcal R_1$, what both codes can say. The loop term acts on holonomy components with range in $\mathcal K_0+\mathcal K_1$, completing the invisible part of each dictionary from what the other code says. Of what neither code can say, only the antisymmetric part $(\nu,-\nu)$ with range in $\mathcal K_0\cap\mathcal K_1$ stays at the prior. When the two row spaces are independent and together have dimension $K$, the invisible flat-truth coordinates are recovered as the loop weight grows, while every tangent loop discrepancy can also be absorbed by an invisible change. Recovery under the flat-truth hypothesis therefore differs from detecting whether that hypothesis holds.

\subsection{Without linearization}

The next results use unrestricted group-valued dictionaries and are exact without linearization. They identify the configurations consistent with grounded messages and exactly flat loops. Restricting dictionaries to a single-exponential image requires intersecting the stated solution sets with that image. The following two statements hold for the grounded model \cref{it1:grounded-model} itself, in the constraint form of the loop term.

\propositionheading{Exact reciprocal identification}{it1:exact-pair}
Let $G=\GL(K,\R)$ and let the true dictionaries of two agents be reciprocal, $\bar\theta^\star_{10}=(\bar\theta^\star_{01})^{-1}$. The pairs $(\bar\theta_{01},\bar\theta_{10})$ that reproduce the message laws of both edges for every referent and satisfy $\bar\theta_{01}\bar\theta_{10}=I$ are exactly
\begin{equation}
\bar\theta_{01}=\bar\theta^\star_{01}s,\qquad\bar\theta_{10}=s^{-1}\bar\theta^\star_{10},\qquad
s\in\operatorname{Stab}(\bar T_1)\cap\operatorname{Stab}(\bar T_0\bar\theta^\star_{01}).
\label{it1:exact-pair-eq}
\end{equation}
The pair is determined uniquely exactly when the stacked map $(\bar T_1;\bar T_0\bar\theta^\star_{01})$ is injective; otherwise the solutions form a set of dimension $K\dim\ker(\bar T_1;\bar T_0\bar\theta^\star_{01})$.
\begin{proof}
By \cref{it1:identifiability}(i), the edge $1\to0$ reproduces its law exactly when $\bar\theta_{01}=\bar\theta^\star_{01}s$ with $s\in\operatorname{Stab}(\bar T_1)$. Reciprocity forces $\bar\theta_{10}=s^{-1}\bar\theta^\star_{10}$, and the edge $0\to1$ reproduces its law exactly when $(\bar\theta^\star_{10})^{-1}s^{-1}\bar\theta^\star_{10}=\bar\theta^\star_{01}s^{-1}(\bar\theta^\star_{01})^{-1}\in\operatorname{Stab}(\bar T_0)$, that is, when $\bar T_0\bar\theta^\star_{01}s^{-1}=\bar T_0\bar\theta^\star_{01}$. The intersection of the two stabilizers is the stabilizer of the stacked map, the set of $s$ with $s-I$ mapping $\R^K$ into its kernel, which is open in an affine space of dimension $K$ times the dimension of that kernel.
\end{proof}

The map $\bar T_0\bar\theta^\star_{01}$ is agent $0$'s code expressed in agent $1$'s reference coordinates, so the pair is identified exactly when, between them, the two agents can say everything.

\propositionheading{Exact flat identification}{it1:exact-flat}
Let $G=\GL(K,\R)$ and $\Gamma$ be connected, and suppose that every configuration whose holonomies on $\mathcal L$ are the identity is a coboundary configuration, $\bar\theta_{ij}=k_ik_j^{-1}$ for some $k\in G^{\mathcal V}$. Let the truth be a coboundary configuration, $\bar\theta^\star_{ij}=g_ig_j^{-1}$, and put $\bar T'_e=\bar T_eg_j$ for $e:j\to i$. The configurations that reproduce every message law and have identity holonomy on $\mathcal L$ are exactly
\begin{equation}
\bar\theta_{ij}=g_iu_iu_j^{-1}g_j^{-1},\qquad u\in G^{\mathcal V},\qquad\bar T'_e(u_i-u_j)=0\text{ for every edge }e:j\to i,
\label{it1:exact-flat-eq}
\end{equation}
and two families $u,u'$ give the same configuration exactly when $u'_v=u_vc$ for all $v$ and one $c\in G$. The truth is the only such configuration exactly when the space $H^0$ of \cref{it1:sheaf-count}, formed with the codes $\bar T'_e$, consists of constant sections. The hypothesis holds for the complete exchange with every two-cycle and every directed triangle.
\begin{proof}
By hypothesis $\bar\theta_{ij}=k_ik_j^{-1}$; write $k_v=g_vu_v$. By \cref{it1:identifiability}(i), the edge $e$ reproduces its law exactly when $(\bar\theta^\star_{ij})^{-1}\bar\theta_{ij}=g_ju_iu_j^{-1}g_j^{-1}\in\operatorname{Stab}(\bar T_e)$, that is, when $\bar T'_eu_iu_j^{-1}=\bar T'_e$, that is, when $\bar T'_eu_i=\bar T'_eu_j$. Two families give the same configuration exactly when $(u_i')^{-1}u_i=(u_j')^{-1}u_j$ on every edge, which on a connected graph means $u'=uc$. If $H^0$ consists of constant sections, every column of $u$ is constant, so $u_v=u_0$ for all $v$ and the configuration is the truth. If $s\in H^0$ is not constant, $u_v=I+ts_vw^T$ with $w\neq0$ and small $t\neq0$ is an invertible family satisfying the condition, and it is not of the form $(c)_v$, so its configuration differs from the truth. For the complete exchange put $k_0=I$ and $k_i=\bar\theta_{i0}$. The two-cycle through $0$ and $i$ gives $\bar\theta_{0i}=k_i^{-1}$, and the directed triangle $0\to j\to i\to0$, whose holonomy is $\bar\theta_{0i}\bar\theta_{ij}\bar\theta_{j0}$, gives $\bar\theta_{ij}=\bar\theta_{0i}^{-1}\bar\theta_{j0}^{-1}=k_ik_j^{-1}$.
\end{proof}

The residual set of \cref{it1:exact-flat-eq} has dimension $K(\dim H^0-K)$ at the truth, the count of \cref{it1:unreachable}. In the noiseless, fully observed case the hypothesis of \cref{it1:exact-flat}, that flat holonomy forces a coboundary configuration, is the synchronizability criterion of group synchronization \citep{GaoBrodzkiMukherjee2021}. What the proposition adds is that each relative element is observed only through a code, so that uniqueness is governed by the sections of the code sheaf.

\remarkheading{A loop family for exact flat identification}{it1:arborescence-loops}
The real cycle-spanning hypothesis of \cref{it1:unreachable} suffices for the tangent count but not for the nonlinear hypothesis of \cref{it1:exact-flat}. On the directed triangle $0\to1\to2\to0$ let $\mathcal L=\{\gamma^2\}$, whose chain $2\gamma$ spans the cycle space, and let all codes have rank one and annihilate $e_3$. With $a^Te_3=-2$, the reflection $s=I+e_3a^T$ satisfies $s^2=I$ and $\bar Ts=\bar T$, so the configuration with dictionary $s$ on one edge and $I$ on the others reproduces every message law of the identity truth, has $H_{\gamma^2}=I$, and has $H_\gamma=s\neq I$; it is not a coboundary configuration. A checkable family exists on every strongly connected graph. Fix a root $r$, an out-arborescence $T_{\mathrm{out}}$ and an in-arborescence $T_{\mathrm{in}}$ at $r$, and for every edge $e:j\to i$ let $\mathcal L$ contain the closed walk from $r$ to $j$ along $T_{\mathrm{out}}$, across $e$, and back to $r$ along $T_{\mathrm{in}}$. Let $k_v$ be the ordered product along the $T_{\mathrm{out}}$ path from $r$ to $v$, with $k_r=I$, and $P_{\mathrm{in}}(v)$ the product along the $T_{\mathrm{in}}$ path from $v$ to $r$. Identity holonomy on the walk of $e$ gives $\bar\theta_{ij}=P_{\mathrm{in}}(i)^{-1}k_j^{-1}$. For the edge of $T_{\mathrm{out}}$ entering $i\neq r$, $k_i=\bar\theta_{ij}k_j$, so $P_{\mathrm{in}}(i)^{-1}=k_i$; this also holds at $r$. Hence $\bar\theta_{ij}=k_ik_j^{-1}$ on every edge, which is the hypothesis of \cref{it1:exact-flat}. For the complete exchange with $r=0$ this family consists of the two-cycles and triangles through agent $0$ used in its proof.

\subsection{Receiver-owned recognition laws}

\remarkheading{Ownership and the loop term}{it1:owned}
Under a variational restriction of the global posterior with factors owned by receivers, each receiver infers its own dictionaries, so its dictionary recognition law is one factor $\mathfrak q_i(\omega_{i\cdot})$ of a law factorized over receivers. A two-cycle couples two receivers and a directed triangle three, so the loop term couples the factors. In the tangent model the factorized law closest to the posterior in $\KL(\mathfrak q\Vert\text{posterior})$ has the exact posterior means and, as covariances, the inverses of the diagonal blocks of the posterior precision \citep[Section~10.1.2]{Bishop2006}. These are smaller in the Loewner order than the exact marginal covariances. For this one Gaussian target, the factorization therefore leaves the estimates of the dictionaries unchanged and makes each receiver overconfident about them. This is the global free energy restricted to a factorized family. The law-level per-agent model of \cref{ch:link-dynamics}, in which each receiver integrates its loop factors against the peers' dictionary laws, is a different objective: its fixed-point means need not equal the posterior means, and its variances differ as well (\cref{it2:echo}). Exact means are a feature of Gaussian targets: for a non-Gaussian posterior a factorized law generally has means different from the exact ones, so for general laws the factorization can change the estimates as well. To update its factor, receiver $i$ needs the current means of the dictionaries that other agents hold along its loops: a loop prior can be used only if agents communicate about how they read each other, in addition to communicating about the world.

\section{Extensions beyond the declared experiments}
\label{it1:ch-next}

\openproblemheading{Learning the dictionaries beyond the tangent model}{it1:open-learning}
\Cref{it1:learning,it1:complementary} give the exact recognition law of the tangent model at every number of exchanges and its limit as $n\to\infty$. Their $\lambda\to\infty$ statements hold for every positive-definite anchor and take $n\to\infty$ first. The two limits do not commute when the truth has detectable dissonance: imposing the flat constraint before the data defines a misspecified model, whose posterior concentrates on the $F_\upsilon$-closest points of $\ker B$ to the truth, with ties broken by the prior \citep{Berk1966}, and distorts the identified coordinates. \Cref{it1:exact-pair,it1:exact-flat} treat the constraint form of the loop term for a flat truth without linearization, under the explicit loop hypothesis of \cref{it1:exact-flat}. The detectable dissonance vanishes exactly when some flat configuration reproduces every message law, so it decides whether flatness can be rejected; the rate at which a test of flatness rejects is governed by the Fisher metric $F_\upsilon$ rather than by the loop metric $W$, and is not computed here. The remaining questions have separate scopes. For the full grounded model \cref{it1:grounded-model} with a prior on $G$ or on a declared exponential chart, establish the limiting conditional law on the identified fiber $\prod_e\bar\theta^\star_e\operatorname{Stab}(\bar T_e)$, or on its preimage in the chart, and determine the selection within the fiber for a dissonant truth and a finite loop weight. The nonlinear covariance-dressed defect and the law-comparison penalty defined above lead to different learning problems; the latter requires its own regularity and stabilizer analysis. Another extension treats latent routing labels, which lower the information below the routed sum of \cref{it1:routed}, and loop weights built from the routing weights. For location-family noise and a non-Gaussian anchor, determine the limit of the recognition law on the fiber, which is no longer a conditioned Gaussian, and whether the filled directions are still recovered exactly as the loop weight grows.

\openproblemheading{Codes as part of the species model}{it1:open-codes}
Let codes be learned as well, so that a shared language is a joint fixed point of code and dictionary learning. The expression-map dynamics of \citet{HansenGhrist2021} is the vector-valued analogue on a sheaf. Determine when the learned codes of a species become related by one frame change, as the shared-language hypothesis of \cref{it1:code} assumes, and when learning drives code-visible dissonance to zero while leaving dissonance in unread directions. When senders choose codes under objectives different from the receiver's, code choice is an information-design problem, of Bayesian-persuasion type when the sender commits to its code \citep{KamenicaGentzkow2011}; the receiver-side results of this chapter then describe one side of a signaling game.

\openproblemheading{Dictionaries in the free-energy pictures}{it1:open-free-energy}
The Gaussian tangent learning calculation fixes a normalized likelihood and an anchored loop prior. Extending it to fully nonlinear receiver-owned dictionaries and learned codes requires separate global, per-agent, and reciprocal objectives, with every law dependence differentiated in the chosen objective. Whether those dynamics reduce code-visible dissonance is a further dynamical question; posterior conditioning within the tangent model alone does not establish it.

\openproblemheading{Beyond Gaussian codes}{it1:open-general}
\Cref{it1:loss,it1:noisy,it1:bandwidth} already hold for general laws and measurable codes, including codes into finite alphabets such as words, and the consensus theorem holds for laws as well as means. \Cref{it1:location} extends identifiability to every location family of noise. Extend the identifiability, learning and dissonance results to codes whose message laws are not a location family, such as codes into finite alphabets, where the Fisher form is replaced by the information in the induced discrete channel and the stabilizer by the subgroup that leaves the channel unchanged, and to agents of different kinds, whose dictionaries are not invertible.

The relation to discourse sheaves is the use of private content spaces and restriction maps into shared message spaces \citep{HansenGhrist2021}. The reconstruction step additionally specifies a receiver conditional law, while invertible dictionaries supply loop products. Unknown frame conventions have also been studied as communication constraints, with frame changes treated as gauge data \citep{BartlettRudolphSpekkens2007,vanEnk2006}. Linear interpretation maps between agents' value spaces, with beliefs lost in their null spaces, are proposed by \citet{Amornbunchornvej2025}, and meaning as maps between individual conceptual spaces, with agreement as a fixed point, by \citet{WarglienGardenfors2013}. Communication as active inference between agents whose shared generative model makes their exchange converge is developed by \citet{FristonFrith2015} and \citet{FristonEtAl2020Language}, and language as efficient compression by \citet{ZaslavskyEtAl2018}. Relative to these sources the present chapter adds frame-relative laws on standard Borel spaces, receiver reconstruction with the exact loss of \cref{it1:loss}, identification of receiver-owned dictionaries modulo code stabilizers, directed loop products with independent opposite dictionaries, and the exact split of a loop prior into filled and detectable directions. These comparisons do not identify graph loops at one context with transport between contexts.
